\documentclass[11pt]{letter}
\usepackage[margin=1in]{geometry}
\usepackage{hyperref}
\hypersetup{colorlinks=true, urlcolor=blue}

\signature{Quang-Vinh Dang}
\address{Quang-Vinh Dang \\ British University Vietnam \\ Hung Yen, Vietnam \\ vinh.dq4@buv.edu.vn}

\begin{document}

\begin{letter}{Editor-in-Chief\\IEEE Transactions on Knowledge and Data Engineering}

\opening{Dear Editor,}

I am pleased to submit the manuscript, ``Fair, Non-Progressive Influence
Control Under Backfire and Parameter Uncertainty,'' for consideration for
publication in IEEE Transactions on Knowledge and Data Engineering.

This paper studies influence maximization under a combined relaxation of
the four assumptions underlying the classical Kempe--Kleinberg--Tardos
guarantee: progressive diffusion, monotone submodularity, known influence
parameters, and the absence of a fairness requirement. Real deployments
routinely violate all four at once -- influence can be actively reversed
through negative word-of-mouth, edge probabilities are estimated rather
than known, and social interventions must not neglect a disadvantaged
subpopulation. We make the following contributions:

\begin{itemize}
\item A sandwich approximation built from two monotone-submodular
  surrogates, giving a $\rho(1-1/e)$ guarantee for every backfire
  intensity $\beta \in [0,1)$, with a directly computable per-instance
  certificate, and a proof that $(1-1/e)$-hardness survives backfire
  exactly.
\item A closed-form Lagrangian index for the resulting 2-state/3-action
  per-node control problem, replacing a bisection-plus-value-iteration
  inner loop with an $O(Km + n\log n)$-per-round procedure -- roughly
  two orders of magnitude faster in our measurements.
\item Identification and correction of a population-weighting pathology
  in an existing isoelastic-welfare fairness objective that
  systematically under-protects small groups.
\item MF-BWI-Fair, a control policy combining these results with
  mean-field belief propagation and Bayesian parameter tracking,
  validated against independently reimplemented published baselines
  (including a direct benchmark of our IMM reimplementation against the
  original authors' own C++ code) on a primary synthetic benchmark, four
  further synthetic graphs, and three real Facebook ego-networks at
  increasing scale.
\end{itemize}

We believe this work is well suited to the scope of IEEE Transactions on
Knowledge and Data Engineering: it connects submodular/RIS-style
influence-maximization theory to restless-bandit sequential control
theory -- a connection we believe is novel for a node-level
activation/maintenance control problem -- while remaining grounded in
classical, guarantee-preserving algorithmic techniques and a thorough,
honestly reported empirical evaluation.

This manuscript is original, has not been published previously, and is
not currently under consideration for publication elsewhere. I have no
conflicts of interest to declare. Thank you for considering our
submission; I look forward to the reviewers' feedback.

\closing{Sincerely,}

\end{letter}
\end{document}
